\documentclass[12pt]{article}

\input{../../_assets/preambulo.tex}


\begin{document}

    % 1. Foto de fondo
    % 2. Título
    % 3. Encabezado Izquierdo
    % 4. Color de fondo
    % 5. Coord x del titulo
    % 6. Coord y del titulo
    % 7. Fecha

    
    \input{../../_assets/portada}
    \portadaExamen{ffccA4.jpg}{Geometría III\\Examen XVIII}{Geometría III. Examen XVIII}{MidnightBlue}{-8}{28}{2026}{Roxana Acedo Parra}

    \begin{description}
        \item[Asignatura] Geometría III.
        \item[Curso Académico] 2025-26.
        \item[Grado] Doble Grado en Ingeniería Informática y Matemáticas.
        \item[Grupo] Único.
        \item[Profesor] Antonio Ros Mulero.
        \item[Descripción] Convocatoria ordinaria.
        \item[Fecha] 19 de enero de 2026.
        \item[Duración] 3 horas.
    \end{description}
    \newpage
    
    \begin{ejercicio}[2,5 puntos] Sean $A, B, C, D$ cuatro puntos de un espacio afín $\mathbb{A}^{3}$ en posición general y $G$ su baricentro.
        \begin{enumerate}
            \item Si consideramos un par de lados opuestos $AB, CD$, probar que los puntos medios de los otros cuatro lados forman un paralelogramo cuyas diagonales se intersecan en $G$.
            \item Sea $f:\mathbb{A}^{3}\rightarrow\mathbb{A}^{3}$ la aplicación afín que satisface
            \[
            f(A)=C, \quad f(B)=D, \quad f(C)=A, \quad f(D)=B
            \]
            Probar que $f$ es una simetría afín y hallar los subespacios que la caracterizan.
        \end{enumerate}
    \end{ejercicio}


    \begin{ejercicio}[2,5 puntos] En el plano euclídeo $\mathbb{R}^{2}$:
        \begin{enumerate}
            \item Determinar explícitamente, respecto del sistema de referencia usual, un movimiento $f$ tal que $f(1,2)=(0,2)$, $f(2,1)=(1,3)$ y $f(0,1)=(1,1)$ y probar que es único.
            \item Clasificar el movimiento y obtener sus elementos geométricos.
        \end{enumerate}
    \end{ejercicio}


    \begin{ejercicio}[2,5 puntos] En $\mathbb{R}^{2}$ y respecto de un sistema de referencia ortonormal, se consideran el punto $Q=(1,0)$, la recta $r\equiv\{x+1=0\}$, un parámetro $m>0$ y el lugar geométrico
    $$\mathcal{C}=\{P\in\mathbb{R}^{2}:d(P,Q)=m~d(P,r)\},$$
    donde $d(P,r)$ denota la distancia del punto $P$ a la recta $r$.
        \begin{enumerate}
            \item Demostrar que $\mathcal{C}$ es una cónica. Clasificarla afínmente y calcular su centro, para todo $m>0$.
            \item Para $m=\sqrt{2}$, determinar los elementos geométricos de la cónica (vértices, ejes, focos, asíntotas, si existen).
        \end{enumerate}
    \end{ejercicio}

    \begin{ejercicio}[2,5 puntos] En el espacio afín $\mathbb{R}^{3}$ y dado un parámetro real $\lambda$, consideramos la cuádrica
    $$\mathcal{Q}\equiv-y^{2}-2xy+2xz+2yz-2\lambda x+\lambda^{2}+\lambda=0.$$
        \begin{enumerate}
            \item Calcular el centro de $\mathcal{Q}$.
            \item Clasificarla afínmente para todos los valores de $\lambda$.
        \end{enumerate}
    \end{ejercicio}

\end{document}